\chapter{Prompt and Context Engineering}
\label{chapter:ContextEngineering}
% EVOLVE-BLOCK-START

A large language model is, at each invocation, a function of its context.\index{Context engineering}
Its instructions, history, domain knowledge, and the results of previous tool calls are all encoded in the same token sequence the model reads.
The model has no persistent state outside this sequence: what it does not read, it does not know.
This constraint makes \emph{context engineering} — the deliberate construction of what enters the context window — the primary lever available to the agent builder.

This chapter distinguishes two related but distinct disciplines.
\emph{Prompt engineering}\index{Prompt engineering} concerns the wording and structure of instructions, especially at the level of a single message or request.
\emph{Context engineering} is broader: it concerns the entire contents of the context window across a multi-step interaction, and the strategies for managing that window as the conversation evolves.
Section~\ref{section:LLMContextFunction} develops the view of the model as a function of its context; Section~\ref{section:PromptEngineering} treats prompt engineering at the level of a single request; Section~\ref{section:ContextEngineeringDetail} scales up to managing the whole window across a multi-step interaction; Section~\ref{section:AutomatedContext} covers the mechanisms that assemble context automatically; and Section~\ref{section:ContextQuality} asks how to tell whether a given context is any good.

\section{The LLM as a Context-Conditioned Function}
\label{section:LLMContextFunction}

Let $V^*$ denote the space of all finite token sequences, and let $\mathcal{A}$ denote the space of possible actions (tool calls plus text responses).
A language model $M$ can be understood informally as a mapping
\begin{equation}
M : V^* \to \Delta(\mathcal{A}) ,
\label{equation:ModelAsFunction}
\end{equation}
where $\Delta(\mathcal{A})$ is the set of probability distributions over actions.
The context — the full token sequence including system prompt, prior turns, and any injected material — determines this distribution completely.

This view has a direct practical consequence.

\begin{principle}
\label{principle:ContextDeterminesOutput}
The context determines the model's output \emph{distribution}: two identical contexts induce the same distribution over next actions.
The realized \emph{sample} depends further on the sampling temperature (introduced below) and can carry additional nondeterminism from production serving, so identical contexts guarantee identical distributions, not identical outputs.
The practical consequence is unchanged: if the agent behaves unexpectedly, look first to the context — either something unexpected is present, or something necessary is absent.
\end{principle}

\subsection{The Context Window}

The \emph{context window}\index{Context window} is the finite-length sequence that $M$ receives at a single invocation.
Window sizes for frontier models in 2025–2026 range from roughly 100,000 to 1,000,000 tokens.
This is large but not unlimited: a single research paper is roughly 10,000 tokens; a large codebase can exceed any window.
These figures are best read as a point on a fast-moving trajectory rather than a fixed ceiling: the first 100,000-token window was itself a frontier result in 2023, an order-of-magnitude jump over the few thousand tokens available shortly before (Anthropic, 2023), and the limit has kept climbing since.
The practical lesson survives each increase --- reason about how to \emph{use} a window rather than assume the next one will be large enough to make the question moot.

Even within the window, not all tokens are equal.
Empirical studies consistently find \emph{recency bias} (later tokens receive more weight) and \emph{primacy bias} (early tokens, particularly system-prompt content, also receive elevated attention), with a \emph{lost-in-the-middle} effect (Liu et al., 2024) where relevant information buried in the center of a long context is underused.

\begin{remark}
The lost-in-the-middle effect implies that simply extending the context window does not proportionally extend the effective working memory of the model.
Early framings of large windows treated raw capacity as a near-substitute for careful selection --- the 2023 case for a 100,000-token window emphasized feeding in whole documents and letting the model synthesize across them (Anthropic, 2023) --- but the effect characterized shortly afterward (Liu et al., 2024) showed that capacity and effective working memory are not the same thing.
Structuring the context so that the most important information appears near the beginning or end — or citing it explicitly in the instruction — is more reliable than relying on a long window to make everything equally accessible.
\end{remark}

\subsection{Temperature and Sampling}

The model samples an action from the distribution $M(\mathcal{C})$.
The \emph{temperature}\index{Temperature} parameter divides the logits by the temperature before the softmax: a temperature of~0 produces deterministic (greedy) output, and higher temperatures flatten the distribution and produce more varied output.

For most agentic tasks, low temperature (0--0.3) is preferred: the agent should follow instructions precisely rather than improvising.
Higher temperatures are appropriate when the goal is to generate diverse candidates — for example, during the proposal step of an evolutionary search (Chapter~\ref{chapter:EvolutionaryFrameworks}).

\section{Prompt Engineering}
\label{section:PromptEngineering}

\emph{Prompt engineering}\index{Prompt engineering} is the practice of writing instructions that reliably elicit the desired model behavior.
It is best understood as a form of specification: the prompt specifies what the model should do, just as source code specifies what a program should do, and clarity matters for the same reasons.

\subsection{Principles of Effective Instructions}

A small set of principles accounts for most of the variance in instruction quality.

\begin{enumerate}
\item \textbf{Be specific about the task.}
Vague instructions produce vague outputs.
``Analyze the data'' is worse than ``Fit a linear model to the data in \code{data.csv} and report the $R^2$ on the held-out test split.''
Every ambiguity in the instruction is a degree of freedom the model fills in on its own.

\item \textbf{Specify the output format.}
Tell the model exactly what to produce: a JSON object with named fields, a Python function with a specific signature, a one-paragraph summary followed by a bullet list.
Explicit format specifications reduce post-processing errors and enable structured pipelines.

\item \textbf{Provide examples.}
\emph{Few-shot prompting}\index{Few-shot prompting} places one or more worked examples in the context before the actual task.
Examples are more reliable than abstract descriptions because they show the model what correct output looks like, not just what it should aim for.

\item \textbf{State constraints explicitly.}
If there are things the model must not do — modify certain files, use certain libraries, produce output longer than a certain length — state them as direct prohibitions rather than hoping the model infers them.

\item \textbf{Decompose complex tasks.}
A single instruction that asks the model to perform many steps is more error-prone than a sequence of instructions that each ask for one step.
Where possible, break multi-part tasks into explicit stages.
\end{enumerate}

\subsection{Chain-of-Thought Prompting}

\emph{Chain-of-thought prompting}\index{Chain-of-thought prompting} (Wei et al., 2022) encourages the model to reason through a task step by step before producing a final answer.
The instruction ``Think through this step by step before giving your final answer'' can substantially improve accuracy on tasks that require multi-step reasoning.
The reasoning trace is not wasted computation: it conditions the final answer on intermediate conclusions, which reduces errors that arise from premature closure.

In the context of the agent loop (Chapter~\ref{chapter:AgentLoop}), the model's scratchpad — the reasoning text before the first tool call — is an instance of chain-of-thought reasoning.

\section{Context Engineering}
\label{section:ContextEngineeringDetail}

Prompt engineering focuses on how instructions are written.
Context engineering focuses on what information is placed in the window, when it is placed there, and how it is organized.

\subsection{The Instruction Stack}

A well-engineered context for an agentic session has a layered structure.

\begin{enumerate}
\item \textbf{System prompt.} The invariant instructions that define the agent's role, tools, behavioral norms, and project-specific conventions.
This content is static across sessions but should be reviewed and updated when the project's needs change.

\item \textbf{Task description.} The specific goal of the current session, provided at the start of the conversation.
This is where the researcher specifies what to accomplish today.

\item \textbf{Injected context.} Relevant files, code snippets, prior results, or retrieved knowledge that the agent needs for the current task.
Unlike the system prompt, this content changes with the task.

\item \textbf{Conversation history.} The sequence of prior observations and actions within the session.
This grows as the session proceeds and must be managed to avoid filling the window with irrelevant history.
\end{enumerate}

\begin{remark}
A common failure mode is treating the system prompt as the only lever.
Researchers who spend hours crafting a perfect system prompt but place nothing relevant in the injected-context layer often find the agent making avoidable mistakes — not because the instructions are unclear, but because the model lacks the specific information it needs to act correctly.
\end{remark}

\subsection{Selective Context Injection}

Not everything that is \emph{relevant} should be placed in the context: filling the window with loosely related material can degrade performance as reliably as providing too little.
The right question is not ``what is relevant to the project?'' but ``what does the agent need to know \emph{right now} to take the next action correctly?''

\begin{principle}
\label{principle:SelectiveContext}
Inject the minimum context required for the current step.
If the agent is fitting a model, inject the data schema and the eval harness interface — not the entire codebase.
If the agent is proposing candidate functions, inject the domain notes and the current best result — not the fitting code.
\end{principle}

\subsection{Long-Context Strategies}

When the necessary context exceeds what can be injected directly, three strategies extend the effective reach.

\begin{enumerate}
\item \textbf{Retrieval.} Index a large corpus and retrieve the most relevant chunks for the current task.
This is the subject of Chapter~\ref{chapter:Retrieval}.
Retrieval and long context are a genuine tradeoff rather than a settled choice: when a query must synthesize evidence spread across an entire document, loading the whole text can beat retrieving fragments of it (Anthropic, 2023), whereas retrieval wins when the relevant knowledge is a small, locatable slice of a corpus far larger than any window.

\item \textbf{Summarization.} Compress long documents or conversation histories into compact summaries before injection.
The summary trades recall for brevity; its quality determines how much information survives.

\item \textbf{Reference, not inclusion.} Give the agent a pointer to a file and the ability to read it, rather than injecting the file's contents.
This works when the agent can selectively read only the sections it needs.
\end{enumerate}

\section{Automated Context Management}
\label{section:AutomatedContext}

\emph{Automated context management}\index{Automated context management} refers to mechanisms that pre-populate the context without requiring the researcher to specify what to inject for every session.
The goal is to make the rig \emph{self-contextualizing}: when the agent opens a project, it automatically knows what it is working on, what tools are available, and what state the project is in.

\subsection{Hooks and Triggers}

A \emph{hook}\index{Hook} is a shell command or script that runs automatically in response to a trigger event — for example, when the agent starts a session, before a tool call, or after a file is modified.
Hooks can inject context by:
\begin{itemize}
\item Running \code{git status} and \code{git log} to summarize recent changes,
\item Loading the results of the most recent eval run,
\item Summarizing which files have changed since the last session.
\end{itemize}

\subsection{Instruction File Hierarchy}

The instruction file hierarchy discussed in Section~\ref{section:InstructionFiles} is itself a form of automated context management: the \code{CLAUDE.md} at the root of the project is loaded automatically at the start of every session.
A well-maintained \code{CLAUDE.md} means the agent starts every session knowing the project's conventions, the location of key files, and the one-command run procedure.

\subsection{The Model Context Protocol}

Where hooks and instruction files inject context with project-specific scripts, the \emph{Model Context Protocol} (MCP)\index{Model Context Protocol} standardizes the mechanism (Anthropic, 2024).
MCP is an open client--server protocol: a \emph{server} exposes tools, resources, and prompts through a typed interface, and any MCP-aware \emph{client} (an agentic CLI, an editor, a chat application) can attach to it and offer its capabilities to the model.
Its value is composability.
A single database server, documentation server, or issue-tracker server, written once, works across every client that speaks the protocol, rather than being re-implemented as a bespoke hook per project.
In the vocabulary of this chapter, MCP is automated context management with a standard interface: it is the protocol-level counterpart to the \code{CLAUDE.md} hierarchy, supplying tools and retrievable resources on demand where the instruction file supplies static conventions.

\section{Measuring Context Quality}
\label{section:ContextQuality}

How does one know whether a context engineering change is an improvement?
The answer is the same as for any other engineering choice: measure it.

\begin{enumerate}
\item \textbf{Ablation.} Remove one element of the context (a retrieved document, a few-shot example, a section of the system prompt) and measure whether performance on a held-out task set changes.
If removing it makes no difference, it is not contributing.

\item \textbf{Prompt sensitivity.} Run the same task with several small variations of the instruction and observe the variance in output.
High variance indicates the model is sensitive to phrasing, which in turn indicates the instruction is under-specified.

\item \textbf{Leaderboard comparison.} In the course, the primary measure of context quality is whether a change improves the agent's score on the shared evaluation harness.
A grounded agent (Challenge~3) beats an ungrounded agent, or it does not — the harness decides.
\end{enumerate}

\section*{Further Reading}

\begin{small}
\begin{enumerate}
\item Wei, J., et al., ``Chain-of-Thought Prompting Elicits Reasoning in Large Language Models,'' \emph{NeurIPS}, 2022 — the empirical case for scratchpad reasoning.
\item Liu, N. F., et al., ``Lost in the Middle: How Language Models Use Long Contexts,'' \emph{TACL}, 2024 — the empirical basis for the lost-in-the-middle effect.
\item Anthropic, \emph{Prompt Engineering Guide} — practical instruction-writing reference for Claude; available at \href{https://docs.anthropic.com/prompt-engineering}{docs.anthropic.com}.
\item Anthropic, \emph{Claude Code: Model Context Protocol} — specification for automated context injection; available in the course workshop materials.
\end{enumerate}
\end{small}
% EVOLVE-BLOCK-END
